\subsection{General Statements}

\begin{lemma}
    Let $T \colon V \to W$ be a linear map of finite dimensional vector spaces and let $A$ be the matrix representation of $T$ with respect to some basis. Then the matrix representation of the dual map $T^\ast \colon W^\ast \to V^\ast$ with respect to the corresponding dual basis is given by $A^T$.
\end{lemma}

\begin{proof}
    Let $X \colon \R^n \to V$ and $Y \colon \R^m \to W$ be linear isomorphisms given by the chosen basis, i.e.

    \[\begin{tikzcd}
	V & W \\
	{\mathbb{R}^n} & {\mathbb{R}^m}
	\arrow["T", from=1-1, to=1-2]
	\arrow["X"', from=2-1, to=1-1]
	\arrow["A", from=2-1, to=2-2]
	\arrow["Y"', from=2-2, to=1-2]
    \end{tikzcd}\]
    commutes. Moreover, let $\phi \colon \R^n \to (\R^n)^\ast$, $\psi \colon \R^m \to (\R^m)^\ast$ be the canonical isomorphisms induced by the standard basis $e_1,\ldots,e_n$ and $d_1,\ldots,d_m$ on $\R^n$ and $\R^m$, respectively. We then get the following commutative diagram
    \[\begin{tikzcd}
	{V^\ast} & {W^\ast} \\
	{(\mathbb{R}^n)^\ast} & {(\mathbb{R}^m)^\ast} \\
	{\mathbb{R}^n} & {\mathbb{R}^m}
	\arrow["{X^\ast}", from=1-1, to=2-1]
	\arrow["{T^\ast}"', from=1-2, to=1-1]
	\arrow["{Y^\ast}", from=1-2, to=2-2]
	\arrow["{A^\ast}"', from=2-2, to=2-1]
	\arrow["\phi"', from=3-1, to=2-1]
	\arrow["\psi"', from=3-2, to=2-2]
	\arrow["B"', from=3-2, to=3-1]
    \end{tikzcd}\]
    where $B$ is the desired matrix representation of $T^\ast$. Write $B = (b_{ij})$. Then $b_{ij}$ is given as the $i$-th component of the vector $Bd_j$. By the above diagram this is precisely
    \[
    b_{ij} = (d_j^\ast \circ A)(e_i) = a_{ji},
    \]
    i.e. $B = A^T$.
\end{proof}

\begin{lemma}
	Let $V, V_1,\ldots,V_k$ be finite dimensional vector spaces. There exist natural isomorphisms
	\begin{align*}
		V_1^\ast \otimes \ldots \otimes V_k^\ast & \xrightarrow{\cong} (V_1 \otimes \ldots \otimes V_k)^\ast \\
		\Lambda^k V^\ast &\xrightarrow{\cong} (\Lambda^k V)^\ast \\
		\Sigma^k V^\ast &\xrightarrow{\cong} (\Sigma^k V)^\ast \\
	\end{align*}
	where $\Lambda^k V$ denotes the $k$-th exterior power and $\Sigma^k V$ the $k$-th symmetric power (see section \ref{linalg:sec:the_exterior_algebra}).
\end{lemma}

\begin{proof}
	Define the first isomorphism $V_1^\ast \otimes \ldots \otimes V_k^\ast \to (V_1 \otimes \ldots \otimes V_k)^\ast $ by 
	\[
	\phi_1 \otimes \ldots \otimes \phi_k \mapsto \left(v_1 \otimes \ldots \otimes v_k \mapsto \prod_{i=1}^k \phi_i(v_i)\right).
	\]
	using the universal property of the tensor product. This map is clearly injective, hence indeed an isomorphism. Moreover this map induces the second isomorphism $\Lambda^k V^\ast \to (\Lambda^k V)^\ast$ by 
		\[\begin{tikzcd}
		{(V^\ast)^{\otimes k}} & {(V^{\otimes k})^\ast} \\
		{\Lambda^kV^\ast} & {(\Lambda^kV)^\ast}
		\arrow["\cong", from=1-1, to=1-2]
		\arrow[from=1-2, to=2-2]
		\arrow[from=2-1, to=1-1]
		\arrow[from=2-1, to=2-2]
		\end{tikzcd}\]
	where the vertical maps are (induced by) the inclusions $\Lambda^kV^\ast \to (V^\ast)^{\otimes k}$ and $\Lambda^kV \to V^{\otimes k}$. To see that this does in fact define an isomorphism note that an inverse is explicitly given by 
		\[\begin{tikzcd}
		{(V^\ast)^{\otimes k}} & {(V^{\otimes k})^\ast} \\
		{\Lambda^kV^\ast} & {(\Lambda^kV)^\ast}
		\arrow["\cong", from=1-1, to=1-2]
		\arrow["{\mathrm{Alt}}"', from=1-1, to=2-1]
		\arrow["{\mathrm{Alt}^\ast}"', from=2-2, to=1-2]
		\arrow[from=2-2, to=2-1]
		\end{tikzcd}\]
	where $\mathrm{Alt}$ denotes the antisymmetrization map (again see section \ref{linalg:sec:the_exterior_algebra}). The third isomorphism $\Sigma^k V^\ast \to (\Sigma^k V)^\ast$ is defined analogously using the symmetrization map instead. Naturality of the defined isomorphisms follows directly from the definiton by a straightforward computation.
\end{proof}